\documentclass[a4paper,10pt]{report}\input{../0.Préambule/Préambule_cours_prof}\begin{document}

\pagestyle{empty}
\newpage
\hspace{8cm}\noindent NOM, Prénom et classe : \dotfill

\vspace{-1cm}
\section*{Parallélisme et angles SUJET A}
\addcontentsline{toc}{section}{Parallélisme et angles SUJET A}

\noindent \textit{Il sera tenu compte dans la notation de la propreté ainsi que de la justification apportée à chacune des réponses.}

\noindent \textit{L'usage de la calculatrice est \textbf{interdite}.} \hfill \textit{Durée : 55 minutes}

\paragraph{Exercice 1}: \textit{Exercice à faire sur une copie séparée .} \hfill \textbf{3 points}

\begin{enumerate}
\item Donner la définition de deux angles adjacents. (On pourra illustrer son propos par une figure)
\item Donner la définition de deux angles supplémentaires %complémentaires
\end{enumerate}

\paragraph{Exercice 2}: \textit{Exercice à faire sur une copie séparée .} \hfill \textbf{9 points}

\medskip
\noindent On considère la figure ci-contre.

\noindent \begin{minipage}{9cm}
\begin{enumerate}
\item Démontre que $(NO)$ et $(LA)$ sont parallèles.
\item Démontre que les angles $\widehat{ALR}$ et $\widehat{NOR}$ ont la même mesure que tu calculeras.
\item Déduis-en la nature du triangle $NOR$.
\end{enumerate}
\end{minipage}
\hspace{5mm}
\begin{minipage}{5cm}
\includegraphics[scale=0.4]{../40.Image/11. Exo-Angles et parallélisme 8.png}
\end{minipage}

\paragraph{Exercice 3}: \textit{Exercice à faire sur cette feuille.} \hfill \textbf{1.5 points}

\medskip
\noindent Sur la figure ci-contre, colorier :

\vspace{-5mm}
\noindent \begin{minipage}{9cm}
\begin{enumerate}
\item en rouge, $2$ angles alternes internes
\item en bleu, $2$ angles supplémentaires
\item en vert, $2$ angles correspondants
\end{enumerate}
\end{minipage}
\hspace{5mm}
\begin{minipage}{5cm}
\begin{center}
\begin{tikzpicture}[line cap=round,line join=round,>=triangle 45,x=1.5cm,y=1.0cm,scale =0.5]
\clip(-1.,0.) rectangle (6.5,6.);
\draw [line width=1.2pt] (-1.,1.)-- (6.,3.);
\draw [line width=1.2pt] (-1.,3.)-- (6.,5.);
\draw [line width=1.2pt] (1.14,5.82)-- (5.12,-0.08);
\begin{scriptsize}
\draw[color=black] (1.26,1.1) node {\large{$(d')$}};
\draw[color=black] (4.42,5.3) node {\large{$(d)$}};
\draw[color=black] (2.2,5.2) node {\large{$(D)$}};
\end{scriptsize}
\end{tikzpicture}
\end{center}
\end{minipage}

\paragraph{Exercice 4}: \textit{Exercice à faire sur cette feuille.} \hfill \textbf{2 points}

\medskip
\noindent Écrire en face de chacune de phrases suivantes « Vrai » ou « Faux ».

\noindent \begin{minipage}{10cm}
\begin{enumerate}
\item $\widehat{xAy}$ et $\widehat{yAz}$ sont adjacents \dotfill \medskip
\item $\widehat{xAy}$ et $\widehat{xAz}$ sont adjacents \dotfill \medskip
\item $\widehat{yAz}$ et $\widehat{zAx}$ sont adjacents \dotfill \medskip
\item $\widehat{xAz}$ et $\widehat{zBt}$ sont adjacents \dotfill \medskip
\end{enumerate}
\end{minipage}
\hspace{5mm}
\begin{minipage}{5cm}
\includegraphics[scale=0.3]{../40.Image/11. DS-Angles et parallélisme 2.png}
\end{minipage}

\paragraph{Exercice 5}: \textit{Exercice à faire sur une copie séparée .} \hfill \textbf{4.5 points}

\medskip
\noindent On considère la figure ci-contre, dans laquelle les droites (DE) et (CB) sont
parallèles.

\noindent \begin{minipage}{11cm}
\begin{enumerate}
\item Calculer la mesure de l'angle $\widehat{ACB}$ . Justifier. \medskip
\item Calculer la mesure de l'angle $\widehat{BEF}$ . Justifier. \medskip
\item Calculer la mesure de l'angle $\widehat{CAB}$ . Justifier.
\end{enumerate}
\end{minipage}
\hspace{5mm}
\begin{minipage}{5cm}
\includegraphics[scale=0.5]{../40.Image/11. DS-Angles et parallélisme 4.png}
\end{minipage}

\newpage
\hspace{8cm}\noindent NOM, Prénom et classe : \dotfill

\vspace{-1cm}
\section*{Parallélisme et angles SUJET B}
\addcontentsline{toc}{section}{Parallélisme et angles SUJET B}

\noindent \textit{Il sera tenu compte dans la notation de la propreté ainsi que de la justification apportée à chacune des réponses.}

\noindent \textit{L'usage de la calculatrice est \textbf{interdite}.} \hfill \textit{Durée : 55 minutes}

\paragraph{Exercice 1}: \textit{Exercice à faire sur une copie séparée .} \hfill \textbf{3 points}

\begin{enumerate}
\item Donner la définition de deux angles opposés par le sommet. (On pourra illustrer son propos par une figure)
\item Donner la définition de deux angles complémentaires
\end{enumerate}

\paragraph{Exercice 2}: \textit{Exercice à faire sur une copie séparée .} \hfill \textbf{9 points}

\medskip
\noindent On considère la figure ci-contre.

\noindent \begin{minipage}{9cm}
\begin{enumerate}
\item Démontre que $(NO)$ et $(LA)$ sont parallèles.
\item Démontre que les angles $\widehat{ALR}$ et $\widehat{NOR}$ ont la même mesure que tu calculeras.
\item Déduis-en la nature du triangle $NOR$.
\end{enumerate}
\end{minipage}
\hspace{5mm}
\begin{minipage}{5cm}
\includegraphics[scale=0.4]{../40.Image/11. Exo-Angles et parallélisme 8.png}
\end{minipage}

\paragraph{Exercice 3}: \textit{Exercice à faire sur cette feuille.} \hfill \textbf{1.5 points}

\medskip
\noindent Sur la figure ci-contre, colorier :

\vspace{-5mm}
\noindent \begin{minipage}{9cm}
\begin{enumerate}
\item en rouge, $2$ angles supplémentaires
\item en bleu, $2$ angles correspondants
\item en vert, $2$ angles alternes internes
\end{enumerate}
\end{minipage}
\hspace{5mm}
\begin{minipage}{5cm}
\begin{center}
\begin{tikzpicture}[line cap=round,line join=round,>=triangle 45,x=1.5cm,y=1.0cm,scale =0.5]
\clip(-1.,0.) rectangle (6.5,6.);
\draw [line width=1.2pt] (-1.,1.)-- (6.,3.);
\draw [line width=1.2pt] (-1.,3.)-- (6.,5.);
\draw [line width=1.2pt] (1.14,5.82)-- (5.12,-0.08);
\begin{scriptsize}
\draw[color=black] (1.26,1.1) node {\large{$(d')$}};
\draw[color=black] (4.42,5.3) node {\large{$(d)$}};
\draw[color=black] (2.2,5.2) node {\large{$(D)$}};
\end{scriptsize}
\end{tikzpicture}
\end{center}
\end{minipage}

\paragraph{Exercice 4}: \textit{Exercice à faire sur cette feuille.} \hfill \textbf{2 points}

\medskip
\noindent Écrire en face de chacune de phrases suivantes « Vrai » ou « Faux ».

\noindent \begin{minipage}{10cm}
\begin{enumerate}
\item $\widehat{xAy}$ et $\widehat{yAz}$ sont adjacents \dotfill \medskip
\item $\widehat{xAy}$ et $\widehat{xAz}$ sont adjacents \dotfill \medskip
\item $\widehat{yAz}$ et $\widehat{zAx}$ sont adjacents \dotfill \medskip
\item $\widehat{xAz}$ et $\widehat{zBt}$ sont adjacents \dotfill \medskip
\end{enumerate}
\end{minipage}
\hspace{5mm}
\begin{minipage}{5cm}
\includegraphics[scale=0.3]{../40.Image/11. DS-Angles et parallélisme 1.png}
\end{minipage}

\paragraph{Exercice 5}: \textit{Exercice à faire sur une copie séparée .} \hfill \textbf{4.5 points}

\medskip
\noindent On considère la figure ci-contre, dans laquelle les droites (DE) et (CB) sont
parallèles.

\noindent \begin{minipage}{11cm}
\begin{enumerate}
\item Calculer la mesure de l'angle $\widehat{ACB}$ . Justifier. \medskip
\item Calculer la mesure de l'angle $\widehat{BEF}$ . Justifier. \medskip
\item Calculer la mesure de l'angle $\widehat{CAB}$ . Justifier.
\end{enumerate}
\end{minipage}
\hspace{5mm}
\begin{minipage}{5cm}
\includegraphics[scale=0.5]{../40.Image/11. DS-Angles et parallélisme 3.png}
\end{minipage}

\end{document}
